% --- Archon physics formalization source begin ---
% archon:physics
% archon:covers PhyXMiniProblems/problem_phyx_mini_0982.lean
% archon:source-report reports/phyx_mini/problem_phyx_mini_0982.source.json

\chapter{Physics problem phyx\_mini\_0982}
\label{ch:PhyXMiniProblems_problem_phyx_mini_0982}

\paragraph{Problem source.}
\#\# Physical scenario

A conducting spherical shell with radius \textbackslash{}( r = 10.0\textbackslash{} \textbackslash{}text\{cm\} \textbackslash{}) spins about a vertical axis twice every second in the presence of a constant magnetic field \textbackslash{}( B = 1.00\textbackslash{} \textbackslash{}text\{T\} \textbackslash{}) that points downward. Two conducting rods supported by a frame contact the sphere with conducting brushes and extend away from the sphere radially. One rod extends from the top of the sphere and the other forms a \textbackslash{}( 60.0\textasciicircum{}\textbackslash{}circ \textbackslash{}) angle with vertical, as shown. The outer ends of the rods are connected to each other by a conducting wire that includes a \textbackslash{}( 10.0\textbackslash{} \textbackslash{}Omega \textbackslash{}) resistor.

\#\# Question

What is the magnitude of the current in the wire?

\#\# Answer choices

- A: A: 0.473mA

- B: B: 4.73mA

- C: C: 47.3mA

- D: D: 0mA

\#\# Figure caption (auxiliary; use the image as primary evidence)

The image depicts a spherical object mounted on a conical stand. The sphere has the following labeled elements:

- A vertical rod is connected to the top of the sphere, with an arrow pointing horizontally to the right, labeled "2 Hz."
- There is another rod attached at an angle of 60 degrees from the vertical rod. This rod has a resistor symbol next to it, labeled "10.0 Ω."
- An arrow pointing downward on the left is labeled "B = 1.00 T," indicating a magnetic field.
- The radius of the sphere is labeled "r = 10.0 cm."

The image likely represents a setup for demonstrating concepts related to rotational motion and electromagnetism.

\#\# Dataset metadata

Electromagnetic Induction; Physical Model Grounding Reasoning, Multi-Formula Reasoning

\paragraph{Recorded answer/context.}
B: B: 4.73mA

\paragraph{Figure/image path.}
/root/proposal\_for\_physic/science-mango/phyx\_mini\_run/phyx\_data/test\_image/982.png

\paragraph{Formalization target.}
create a compiling Lean file with sorry bodies at `PhyXMiniProblems/problem\_phyx\_mini\_0982.lean`. The Lean declarations must preserve the physical quantities, dimensions or dimensional roles, figure labels, governing-law hypotheses, and final relation expressed by this problem.
Use Mathlib/Physlib names found through LeanExplore where available. If a physics API is missing, introduce faithful local abstractions rather than scalar placeholder aliases.

\begin{theorem}[Physics formalization target]
\label{thm:physics:phyx_mini_0982:target}
\lean{PhyXMiniProblems.ProblemPhyXMini0982.rotating_sphere_wire_current}
\uses{def:physics:phyx-mini-0982:phyxminiproblems-problemphyxmini0982-currentinamperes, def:physics:phyx-mini-0982:phyxminiproblems-problemphyxmini0982-rotatingspherecircuitsetup, def:physics:phyx-mini-0982:phyxminiproblems-problemphyxmini0982-matchesconductingspherescenario, def:physics:phyx-mini-0982:phyxminiproblems-problemphyxmini0982-matchesprimaryrotatingspherefigure, def:physics:phyx-mini-0982:phyxminiproblems-problemphyxmini0982-hasphysicalrotatingsphereparameters, def:physics:phyx-mini-0982:phyxminiproblems-problemphyxmini0982-satisfiesrotationfrequencylaw, def:physics:phyx-mini-0982:phyxminiproblems-problemphyxmini0982-satisfiesrotatingspheremotionalemflaw, def:physics:phyx-mini-0982:phyxminiproblems-problemphyxmini0982-satisfiesexternalcircuitohmslaw, def:physics:phyx-mini-0982:phyxminiproblems-problemphyxmini0982-isuniqueclosestcurrentchoice}
The assigned autoformalize agent should translate this physics problem into Lean declarations in the covered file, with theorem and lemma proof bodies written as `by sorry`.
\end{theorem}
\begin{proof}
This is an autoformalization task, not a proof task. Produce statements that can later be proved by the physics prover without weakening the physical model.
\end{proof}

% --- Archon declaration topology begin ---
\section{Declaration topology}
\begin{definition}[inverse Time Dimension]
  \label{def:physics:phyx-mini-0982:phyxminiproblems-problemphyxmini0982-inversetimedimension}
  \lean{PhyXMiniProblems.ProblemPhyXMini0982.inverseTimeDimension}
  Frequency and angular speed both have inverse-time dimension
\end{definition}
\begin{proof}
  This is the declared model definition of inverse Time Dimension.
\end{proof}
\begin{definition}[magnetic Flux Density Dimension]
  \label{def:physics:phyx-mini-0982:phyxminiproblems-problemphyxmini0982-magneticfluxdensitydimension}
  \lean{PhyXMiniProblems.ProblemPhyXMini0982.magneticFluxDensityDimension}
  Magnetic flux density has SI dimension M T⁻¹ C⁻¹ (tesla)
\end{definition}
\begin{proof}
  This is the declared model definition of magnetic Flux Density Dimension.
\end{proof}
\begin{definition}[electric Potential Dimension]
  \label{def:physics:phyx-mini-0982:phyxminiproblems-problemphyxmini0982-electricpotentialdimension}
  \lean{PhyXMiniProblems.ProblemPhyXMini0982.electricPotentialDimension}
  Electric potential has SI dimension M L² T⁻² C⁻¹ (volt)
\end{definition}
\begin{proof}
  This is the declared model definition of electric Potential Dimension.
\end{proof}
\begin{definition}[electric Current Dimension]
  \label{def:physics:phyx-mini-0982:phyxminiproblems-problemphyxmini0982-electriccurrentdimension}
  \lean{PhyXMiniProblems.ProblemPhyXMini0982.electricCurrentDimension}
  Electric current has SI dimension C T⁻¹ (ampere)
\end{definition}
\begin{proof}
  This is the declared model definition of electric Current Dimension.
\end{proof}
\begin{definition}[electrical Resistance Dimension]
  \label{def:physics:phyx-mini-0982:phyxminiproblems-problemphyxmini0982-electricalresistancedimension}
  \lean{PhyXMiniProblems.ProblemPhyXMini0982.electricalResistanceDimension}
  Electrical resistance has SI dimension M L² T⁻¹ C⁻² (ohm)
\end{definition}
\begin{proof}
  This is the declared model definition of electrical Resistance Dimension.
\end{proof}
\begin{definition}[Length Magnitude]
  \label{def:physics:phyx-mini-0982:phyxminiproblems-problemphyxmini0982-lengthmagnitude}
  \lean{PhyXMiniProblems.ProblemPhyXMini0982.LengthMagnitude}
  A nonnegative, unit-independent physical length
\end{definition}
\begin{proof}
  This is the declared model definition of Length Magnitude.
\end{proof}
\begin{definition}[Frequency Magnitude]
  \label{def:physics:phyx-mini-0982:phyxminiproblems-problemphyxmini0982-frequencymagnitude}
  \lean{PhyXMiniProblems.ProblemPhyXMini0982.FrequencyMagnitude}
  \uses{def:physics:phyx-mini-0982:phyxminiproblems-problemphyxmini0982-inversetimedimension}
  A nonnegative, unit-independent rotation frequency
\end{definition}
\begin{proof}
  This is the declared model definition of Frequency Magnitude.
\end{proof}
\begin{definition}[Angular Speed Magnitude]
  \label{def:physics:phyx-mini-0982:phyxminiproblems-problemphyxmini0982-angularspeedmagnitude}
  \lean{PhyXMiniProblems.ProblemPhyXMini0982.AngularSpeedMagnitude}
  \uses{def:physics:phyx-mini-0982:phyxminiproblems-problemphyxmini0982-inversetimedimension}
  A nonnegative, unit-independent angular speed
\end{definition}
\begin{proof}
  This is the declared model definition of Angular Speed Magnitude.
\end{proof}
\begin{definition}[Magnetic Flux Density Magnitude]
  \label{def:physics:phyx-mini-0982:phyxminiproblems-problemphyxmini0982-magneticfluxdensitymagnitude}
  \lean{PhyXMiniProblems.ProblemPhyXMini0982.MagneticFluxDensityMagnitude}
  \uses{def:physics:phyx-mini-0982:phyxminiproblems-problemphyxmini0982-magneticfluxdensitydimension}
  A nonnegative, unit-independent magnetic-flux-density magnitude
\end{definition}
\begin{proof}
  This is the declared model definition of Magnetic Flux Density Magnitude.
\end{proof}
\begin{definition}[Resistance Magnitude]
  \label{def:physics:phyx-mini-0982:phyxminiproblems-problemphyxmini0982-resistancemagnitude}
  \lean{PhyXMiniProblems.ProblemPhyXMini0982.ResistanceMagnitude}
  \uses{def:physics:phyx-mini-0982:phyxminiproblems-problemphyxmini0982-electricalresistancedimension}
  A nonnegative, unit-independent electrical resistance
\end{definition}
\begin{proof}
  This is the declared model definition of Resistance Magnitude.
\end{proof}
\begin{definition}[Electromotive Force Magnitude]
  \label{def:physics:phyx-mini-0982:phyxminiproblems-problemphyxmini0982-electromotiveforcemagnitude}
  \lean{PhyXMiniProblems.ProblemPhyXMini0982.ElectromotiveForceMagnitude}
  \uses{def:physics:phyx-mini-0982:phyxminiproblems-problemphyxmini0982-electricpotentialdimension}
  A nonnegative, unit-independent electromotive-force magnitude
\end{definition}
\begin{proof}
  This is the declared model definition of Electromotive Force Magnitude.
\end{proof}
\begin{definition}[Electric Current Magnitude]
  \label{def:physics:phyx-mini-0982:phyxminiproblems-problemphyxmini0982-electriccurrentmagnitude}
  \lean{PhyXMiniProblems.ProblemPhyXMini0982.ElectricCurrentMagnitude}
  \uses{def:physics:phyx-mini-0982:phyxminiproblems-problemphyxmini0982-electriccurrentdimension}
  A nonnegative, unit-independent electric-current magnitude
\end{definition}
\begin{proof}
  This is the declared model definition of Electric Current Magnitude.
\end{proof}
\begin{definition}[coherent SIReadout]
  \label{def:physics:phyx-mini-0982:phyxminiproblems-problemphyxmini0982-coherentsireadout}
  \lean{PhyXMiniProblems.ProblemPhyXMini0982.coherentSIReadout}
  Coherent-SI scalar readout of a nonnegative dimensionful magnitude
\end{definition}
\begin{proof}
  This is the declared model definition of coherent SIReadout.
\end{proof}
\begin{definition}[length In Meters]
  \label{def:physics:phyx-mini-0982:phyxminiproblems-problemphyxmini0982-lengthinmeters}
  \lean{PhyXMiniProblems.ProblemPhyXMini0982.lengthInMeters}
  \uses{def:physics:phyx-mini-0982:phyxminiproblems-problemphyxmini0982-lengthmagnitude, def:physics:phyx-mini-0982:phyxminiproblems-problemphyxmini0982-coherentsireadout}
  Read a length magnitude in metres
\end{definition}
\begin{proof}
  This is the declared model definition of length In Meters.
\end{proof}
\begin{definition}[frequency In Hertz]
  \label{def:physics:phyx-mini-0982:phyxminiproblems-problemphyxmini0982-frequencyinhertz}
  \lean{PhyXMiniProblems.ProblemPhyXMini0982.frequencyInHertz}
  \uses{def:physics:phyx-mini-0982:phyxminiproblems-problemphyxmini0982-frequencymagnitude, def:physics:phyx-mini-0982:phyxminiproblems-problemphyxmini0982-coherentsireadout}
  Read a rotation frequency in hertz
\end{definition}
\begin{proof}
  This is the declared model definition of frequency In Hertz.
\end{proof}
\begin{definition}[angular Speed In Radians Per Second]
  \label{def:physics:phyx-mini-0982:phyxminiproblems-problemphyxmini0982-angularspeedinradianspersecond}
  \lean{PhyXMiniProblems.ProblemPhyXMini0982.angularSpeedInRadiansPerSecond}
  \uses{def:physics:phyx-mini-0982:phyxminiproblems-problemphyxmini0982-angularspeedmagnitude, def:physics:phyx-mini-0982:phyxminiproblems-problemphyxmini0982-coherentsireadout}
  Read an angular speed in radians per second
\end{definition}
\begin{proof}
  This is the declared model definition of angular Speed In Radians Per Second.
\end{proof}
\begin{definition}[magnetic Flux Density In Teslas]
  \label{def:physics:phyx-mini-0982:phyxminiproblems-problemphyxmini0982-magneticfluxdensityinteslas}
  \lean{PhyXMiniProblems.ProblemPhyXMini0982.magneticFluxDensityInTeslas}
  \uses{def:physics:phyx-mini-0982:phyxminiproblems-problemphyxmini0982-magneticfluxdensitymagnitude, def:physics:phyx-mini-0982:phyxminiproblems-problemphyxmini0982-coherentsireadout}
  Read a magnetic-flux-density magnitude in teslas
\end{definition}
\begin{proof}
  This is the declared model definition of magnetic Flux Density In Teslas.
\end{proof}
\begin{definition}[resistance In Ohms]
  \label{def:physics:phyx-mini-0982:phyxminiproblems-problemphyxmini0982-resistanceinohms}
  \lean{PhyXMiniProblems.ProblemPhyXMini0982.resistanceInOhms}
  \uses{def:physics:phyx-mini-0982:phyxminiproblems-problemphyxmini0982-resistancemagnitude, def:physics:phyx-mini-0982:phyxminiproblems-problemphyxmini0982-coherentsireadout}
  Read an electrical resistance in ohms
\end{definition}
\begin{proof}
  This is the declared model definition of resistance In Ohms.
\end{proof}
\begin{definition}[electromotive Force In Volts]
  \label{def:physics:phyx-mini-0982:phyxminiproblems-problemphyxmini0982-electromotiveforceinvolts}
  \lean{PhyXMiniProblems.ProblemPhyXMini0982.electromotiveForceInVolts}
  \uses{def:physics:phyx-mini-0982:phyxminiproblems-problemphyxmini0982-electromotiveforcemagnitude, def:physics:phyx-mini-0982:phyxminiproblems-problemphyxmini0982-coherentsireadout}
  Read an electromotive-force magnitude in volts
\end{definition}
\begin{proof}
  This is the declared model definition of electromotive Force In Volts.
\end{proof}
\begin{definition}[current In Amperes]
  \label{def:physics:phyx-mini-0982:phyxminiproblems-problemphyxmini0982-currentinamperes}
  \lean{PhyXMiniProblems.ProblemPhyXMini0982.currentInAmperes}
  \uses{def:physics:phyx-mini-0982:phyxminiproblems-problemphyxmini0982-electriccurrentmagnitude, def:physics:phyx-mini-0982:phyxminiproblems-problemphyxmini0982-coherentsireadout}
  Read an electric-current magnitude in amperes
\end{definition}
\begin{proof}
  This is the declared model definition of current In Amperes.
\end{proof}
\begin{definition}[radians From Degrees]
  \label{def:physics:phyx-mini-0982:phyxminiproblems-problemphyxmini0982-radiansfromdegrees}
  \lean{PhyXMiniProblems.ProblemPhyXMini0982.radiansFromDegrees}
  Convert a degree readout into its dimensionless radian value
\end{definition}
\begin{proof}
  This is the declared model definition of radians From Degrees.
\end{proof}
\begin{definition}[Sphere Contact]
  \label{def:physics:phyx-mini-0982:phyxminiproblems-problemphyxmini0982-spherecontact}
  \lean{PhyXMiniProblems.ProblemPhyXMini0982.SphereContact}
  The two brush-contact locations on the spherical shell
\end{definition}
\begin{proof}
  This is the declared model definition of Sphere Contact.
\end{proof}
\begin{definition}[Axis Orientation]
  \label{def:physics:phyx-mini-0982:phyxminiproblems-problemphyxmini0982-axisorientation}
  \lean{PhyXMiniProblems.ProblemPhyXMini0982.AxisOrientation}
  Orientations needed to state the axis geometry
\end{definition}
\begin{proof}
  This is the declared model definition of Axis Orientation.
\end{proof}
\begin{definition}[Axial Direction]
  \label{def:physics:phyx-mini-0982:phyxminiproblems-problemphyxmini0982-axialdirection}
  \lean{PhyXMiniProblems.ProblemPhyXMini0982.AxialDirection}
  Directions along the vertical rotation axis
\end{definition}
\begin{proof}
  This is the declared model definition of Axial Direction.
\end{proof}
\begin{definition}[Magnetic Field Regime]
  \label{def:physics:phyx-mini-0982:phyxminiproblems-problemphyxmini0982-magneticfieldregime}
  \lean{PhyXMiniProblems.ProblemPhyXMini0982.MagneticFieldRegime}
  Qualitative models for the applied magnetic field
\end{definition}
\begin{proof}
  This is the declared model definition of Magnetic Field Regime.
\end{proof}
\begin{definition}[Conductor Motion]
  \label{def:physics:phyx-mini-0982:phyxminiproblems-problemphyxmini0982-conductormotion}
  \lean{PhyXMiniProblems.ProblemPhyXMini0982.ConductorMotion}
  Motion of a conductor relative to the laboratory frame
\end{definition}
\begin{proof}
  This is the declared model definition of Conductor Motion.
\end{proof}
\begin{definition}[Rotating Sphere Circuit Figure]
  \label{def:physics:phyx-mini-0982:phyxminiproblems-problemphyxmini0982-rotatingspherecircuitfigure}
  \lean{PhyXMiniProblems.ProblemPhyXMini0982.RotatingSphereCircuitFigure}
  \uses{def:physics:phyx-mini-0982:phyxminiproblems-problemphyxmini0982-spherecontact}
  Literal presentation data visible in the primary raster. Numerical fields are the printed scalar labels; the calibration predicate below ties them to dimensionful physical magnitudes
\end{definition}
\begin{proof}
  This is the declared model definition of Rotating Sphere Circuit Figure.
\end{proof}
\begin{definition}[Rotating Sphere Circuit Setup]
  \label{def:physics:phyx-mini-0982:phyxminiproblems-problemphyxmini0982-rotatingspherecircuitsetup}
  \lean{PhyXMiniProblems.ProblemPhyXMini0982.RotatingSphereCircuitSetup}
  \uses{def:physics:phyx-mini-0982:phyxminiproblems-problemphyxmini0982-lengthmagnitude, def:physics:phyx-mini-0982:phyxminiproblems-problemphyxmini0982-frequencymagnitude, def:physics:phyx-mini-0982:phyxminiproblems-problemphyxmini0982-angularspeedmagnitude, def:physics:phyx-mini-0982:phyxminiproblems-problemphyxmini0982-magneticfluxdensitymagnitude, def:physics:phyx-mini-0982:phyxminiproblems-problemphyxmini0982-resistancemagnitude, def:physics:phyx-mini-0982:phyxminiproblems-problemphyxmini0982-electromotiveforcemagnitude, def:physics:phyx-mini-0982:phyxminiproblems-problemphyxmini0982-electriccurrentmagnitude, def:physics:phyx-mini-0982:phyxminiproblems-problemphyxmini0982-spherecontact, def:physics:phyx-mini-0982:phyxminiproblems-problemphyxmini0982-axisorientation, def:physics:phyx-mini-0982:phyxminiproblems-problemphyxmini0982-axialdirection, def:physics:phyx-mini-0982:phyxminiproblems-problemphyxmini0982-magneticfieldregime, def:physics:phyx-mini-0982:phyxminiproblems-problemphyxmini0982-conductormotion, def:physics:phyx-mini-0982:phyxminiproblems-problemphyxmini0982-rotatingspherecircuitfigure}
  Independent physical objects and observables in the experiment. In particular, motionalEmfMagnitude and wireCurrentMagnitude are not defined by the target formula or an answer choice
\end{definition}
\begin{proof}
  This is the declared model definition of Rotating Sphere Circuit Setup.
\end{proof}
\begin{definition}[Matches Conducting Sphere Scenario]
  \label{def:physics:phyx-mini-0982:phyxminiproblems-problemphyxmini0982-matchesconductingspherescenario}
  \lean{PhyXMiniProblems.ProblemPhyXMini0982.MatchesConductingSphereScenario}
  \uses{def:physics:phyx-mini-0982:phyxminiproblems-problemphyxmini0982-rotatingspherecircuitsetup}
  Qualitative physical setup specified by the prose
\end{definition}
\begin{proof}
  This is the declared model definition of Matches Conducting Sphere Scenario.
\end{proof}
\begin{definition}[Matches Primary Rotating Sphere Figure]
  \label{def:physics:phyx-mini-0982:phyxminiproblems-problemphyxmini0982-matchesprimaryrotatingspherefigure}
  \lean{PhyXMiniProblems.ProblemPhyXMini0982.MatchesPrimaryRotatingSphereFigure}
  \uses{def:physics:phyx-mini-0982:phyxminiproblems-problemphyxmini0982-lengthinmeters, def:physics:phyx-mini-0982:phyxminiproblems-problemphyxmini0982-frequencyinhertz, def:physics:phyx-mini-0982:phyxminiproblems-problemphyxmini0982-magneticfluxdensityinteslas, def:physics:phyx-mini-0982:phyxminiproblems-problemphyxmini0982-resistanceinohms, def:physics:phyx-mini-0982:phyxminiproblems-problemphyxmini0982-radiansfromdegrees, def:physics:phyx-mini-0982:phyxminiproblems-problemphyxmini0982-rotatingspherecircuitsetup}
  Primary-raster evidence and calibrations for every printed numerical label. The polar angles are figure-derived geometry, not conclusions about current
\end{definition}
\begin{proof}
  This is the declared model definition of Matches Primary Rotating Sphere Figure.
\end{proof}
\begin{definition}[Has Physical Rotating Sphere Parameters]
  \label{def:physics:phyx-mini-0982:phyxminiproblems-problemphyxmini0982-hasphysicalrotatingsphereparameters}
  \lean{PhyXMiniProblems.ProblemPhyXMini0982.HasPhysicalRotatingSphereParameters}
  \uses{def:physics:phyx-mini-0982:phyxminiproblems-problemphyxmini0982-lengthinmeters, def:physics:phyx-mini-0982:phyxminiproblems-problemphyxmini0982-frequencyinhertz, def:physics:phyx-mini-0982:phyxminiproblems-problemphyxmini0982-magneticfluxdensityinteslas, def:physics:phyx-mini-0982:phyxminiproblems-problemphyxmini0982-resistanceinohms, def:physics:phyx-mini-0982:phyxminiproblems-problemphyxmini0982-rotatingspherecircuitsetup}
  Positivity conditions selecting the nondegenerate physical setup
\end{definition}
\begin{proof}
  This is the declared model definition of Has Physical Rotating Sphere Parameters.
\end{proof}
\begin{definition}[Satisfies Rotation Frequency Law]
  \label{def:physics:phyx-mini-0982:phyxminiproblems-problemphyxmini0982-satisfiesrotationfrequencylaw}
  \lean{PhyXMiniProblems.ProblemPhyXMini0982.SatisfiesRotationFrequencyLaw}
  \uses{def:physics:phyx-mini-0982:phyxminiproblems-problemphyxmini0982-frequencyinhertz, def:physics:phyx-mini-0982:phyxminiproblems-problemphyxmini0982-angularspeedinradianspersecond, def:physics:phyx-mini-0982:phyxminiproblems-problemphyxmini0982-rotatingspherecircuitsetup}
  Conversion from rotation frequency to angular speed: ω = 2πf
\end{definition}
\begin{proof}
  This is the declared model definition of Satisfies Rotation Frequency Law.
\end{proof}
\begin{definition}[Satisfies Rotating Sphere Motional Emf Law]
  \label{def:physics:phyx-mini-0982:phyxminiproblems-problemphyxmini0982-satisfiesrotatingspheremotionalemflaw}
  \lean{PhyXMiniProblems.ProblemPhyXMini0982.SatisfiesRotatingSphereMotionalEmfLaw}
  \uses{def:physics:phyx-mini-0982:phyxminiproblems-problemphyxmini0982-lengthinmeters, def:physics:phyx-mini-0982:phyxminiproblems-problemphyxmini0982-angularspeedinradianspersecond, def:physics:phyx-mini-0982:phyxminiproblems-problemphyxmini0982-magneticfluxdensityinteslas, def:physics:phyx-mini-0982:phyxminiproblems-problemphyxmini0982-electromotiveforceinvolts, def:physics:phyx-mini-0982:phyxminiproblems-problemphyxmini0982-rotatingspherecircuitsetup}
  For an axially magnetized conducting sphere rotating about that axis, the motional-emf magnitude between polar angles θ₁ and θ₂ is (1/2) B ω r² |sin² θ₂ - sin² θ₁|. This relates independent observables and does not state the requested current or its numerical value
\end{definition}
\begin{proof}
  This is the declared model definition of Satisfies Rotating Sphere Motional Emf Law.
\end{proof}
\begin{definition}[Satisfies External Circuit Ohms Law]
  \label{def:physics:phyx-mini-0982:phyxminiproblems-problemphyxmini0982-satisfiesexternalcircuitohmslaw}
  \lean{PhyXMiniProblems.ProblemPhyXMini0982.SatisfiesExternalCircuitOhmsLaw}
  \uses{def:physics:phyx-mini-0982:phyxminiproblems-problemphyxmini0982-resistanceinohms, def:physics:phyx-mini-0982:phyxminiproblems-problemphyxmini0982-electromotiveforceinvolts, def:physics:phyx-mini-0982:phyxminiproblems-problemphyxmini0982-currentinamperes, def:physics:phyx-mini-0982:phyxminiproblems-problemphyxmini0982-rotatingspherecircuitsetup}
  Ohm's law for the current magnitude in the closed external circuit
\end{definition}
\begin{proof}
  This is the declared model definition of Satisfies External Circuit Ohms Law.
\end{proof}
\begin{definition}[Current Answer Choice]
  \label{def:physics:phyx-mini-0982:phyxminiproblems-problemphyxmini0982-currentanswerchoice}
  \lean{PhyXMiniProblems.ProblemPhyXMini0982.CurrentAnswerChoice}
  The four current choices printed with the problem
\end{definition}
\begin{proof}
  This is the declared model definition of Current Answer Choice.
\end{proof}
\begin{definition}[answer Choice In Milliamperes]
  \label{def:physics:phyx-mini-0982:phyxminiproblems-problemphyxmini0982-answerchoiceinmilliamperes}
  \lean{PhyXMiniProblems.ProblemPhyXMini0982.answerChoiceInMilliamperes}
  \uses{def:physics:phyx-mini-0982:phyxminiproblems-problemphyxmini0982-currentanswerchoice}
  Printed magnitude of each answer choice, in milliamperes
\end{definition}
\begin{proof}
  This is the declared model definition of answer Choice In Milliamperes.
\end{proof}
\begin{definition}[current In Milliamperes]
  \label{def:physics:phyx-mini-0982:phyxminiproblems-problemphyxmini0982-currentinmilliamperes}
  \lean{PhyXMiniProblems.ProblemPhyXMini0982.currentInMilliamperes}
  \uses{def:physics:phyx-mini-0982:phyxminiproblems-problemphyxmini0982-electriccurrentmagnitude, def:physics:phyx-mini-0982:phyxminiproblems-problemphyxmini0982-currentinamperes}
  Coherent-SI current readout converted from amperes to milliamperes
\end{definition}
\begin{proof}
  This is the declared model definition of current In Milliamperes.
\end{proof}
\begin{definition}[Is Unique Closest Current Choice]
  \label{def:physics:phyx-mini-0982:phyxminiproblems-problemphyxmini0982-isuniqueclosestcurrentchoice}
  \lean{PhyXMiniProblems.ProblemPhyXMini0982.IsUniqueClosestCurrentChoice}
  \uses{def:physics:phyx-mini-0982:phyxminiproblems-problemphyxmini0982-rotatingspherecircuitsetup, def:physics:phyx-mini-0982:phyxminiproblems-problemphyxmini0982-currentanswerchoice, def:physics:phyx-mini-0982:phyxminiproblems-problemphyxmini0982-answerchoiceinmilliamperes, def:physics:phyx-mini-0982:phyxminiproblems-problemphyxmini0982-currentinmilliamperes}
  A listed choice is uniquely closest to the physical current readout. This definition does not privilege any particular choice or prescribe the current
\end{definition}
\begin{proof}
  This is the declared model definition of Is Unique Closest Current Choice.
\end{proof}
% --- Archon declaration topology end ---
% --- Archon physics formalization source end ---
